\section{An all-degree consequence of quadratic elimination}
\label{sec:constantfibres}

The estimate used for the repeated lines gives a useful result
in arbitrary dimension. It also records the rational marked periods that arise before pole
cancellation. This appendix is not needed for the dimension--degree
boundary.

\begin{theorem}\label{thm:constantfibres}
Let $u,Y_1,\ldots,Y_s$ be independent standard Gaussian variables,
$s\ge1$, and
\[
 P=Y^TAY+B(u)^TY+g(u),\qquad A=A^T,
\]
where $B,g$ have arbitrary degrees. Put $D=\det(I-2tA)$ and
$C=(I-2tA)^{-1}$. Purity is equivalent to the following finite
rational identities:
\begin{align}
 \deg g\le2,\qquad &\deg_u(B^TCB)\le2,\label{eq:fibredegree}\\
 D K=1,\qquad &\gamma+\beta^2/(2K)=0,\label{eq:fibreunit}
\end{align}
where $\alpha u^2+\beta u+\gamma=tg+(t^2/2)B^TCB$ and
$K=1-2\alpha$.

For every pure member, every original marked period is rational,
with possible finite poles only at zeros of $D$. If every
nonzero eigenvalue of $A$ is semisimple, all those poles cancel:
$P$ is safe. More explicitly, with
$b=\max(1,\deg_u B)$, interpreting $b=1$ if $B=0$,
\begin{equation}\label{eq:fibrecutoff}
 \deg_t\Phi_Q\le\bigl(s+(s+1)b\bigr)\deg Q.
\end{equation}
No restriction is imposed on the nilpotent zero-eigenspace.
The bound is sufficient, not asserted optimal.
\end{theorem}
\begin{proof}
Lemma~\ref{lem:quadraticfibrehighdegree} forces the effective
phase to be quadratic. Its higher-degree coefficients say
$tg_j+(t^2/2)[u^j]B^TCB=0$. At $t=0$ these first give $g_j=0$
for $j\ge3$, then the second identity in
\eqref{eq:fibredegree}. The unit series is exactly
\[
 \Phi_1=(DK)^{-1/2}\exp\!\left(\gamma+\frac{\beta^2}{2K}\right).
\]
Squaring the purity identity gives
$DK=\exp(2\gamma+\beta^2/K)$. A rational derivative has zero
residues, whereas the logarithmic derivative of a rational
function has its zero and pole orders as residues. Thus $DK$
is constant, and then so is $2\gamma+\beta^2/K$.
Normalization at zero gives
\eqref{eq:fibreunit}. Conversely those identities give
$\Phi_1=1$.

Set $\mu=D\beta$. It is polynomial because the only denominator
in $C$ is $D$. Under the normalized tilted functional the base
variable $u$ has mean $\mu$ and variance $D$. Conditional on
$u$, the $Y$ mean is $tCB(u)$ and its covariance is $C$.
Wick's rule therefore expresses every original marked period
as a rational function whose denominator is a power of $D$.
The unit period is one, so no additional exponential factor
remains. In particular their span over $\C(t)$ has rank one;
this assertion alone does not prove safety.

Assume now that all nonzero eigenvalues are semisimple. The
generalized eigenspaces of the symmetric matrix are orthogonal
and nondegenerate. Write $A_0$ for its nilpotent zero block.
Its resolvent $C_0$ is polynomial. Partial fractions in $t$
separate its contribution from the simple-pole terms at
nonzero eigenvalues. Thus each corresponding quadratic norm
in $B(u)$ has degree at most two in $u$.

Equation $K=1/D$ forces every nonzero eigenvalue to have
multiplicity one: the left side has at most a simple pole
there, while $1/D$ has its full eigenvalue multiplicity.
For such an eigenvalue $\lambda$, put
\[
 l_\lambda=1-2\lambda t,\quad t_\lambda=(2\lambda)^{-1},
 \quad D_\lambda=D/l_\lambda.
\]
Its scalar source is affine, $B_\lambda(u)=b_\lambda u+c_\lambda$,
because its square has degree at most two. Evaluating $DK=1$
at $t_\lambda$ gives the exact certificate
\begin{equation}\label{eq:fibrepolecertificate}
 t_\lambda^2b_\lambda^2D_\lambda(t_\lambda)=-1,
 \qquad
 \mu(t_\lambda)=-c_\lambda/b_\lambda.
\end{equation}
The second identity follows by evaluating $D\beta$ there.

The joint tilted functional of $u$ and all these nonzero-eigenvalue
coordinates is Gaussian. Its means and covariances are
\begin{align*}
 \E_tu&=\mu,&\E_tY_\lambda
   &=t(b_\lambda\mu+c_\lambda)/l_\lambda,\\
 \operatorname{Var}_t(u)&=D,&
 \operatorname{Cov}_t(u,Y_\lambda)&=tDb_\lambda/l_\lambda,\\
 \operatorname{Cov}_t(Y_\lambda,Y_\nu)
   &=t^2b_\lambda b_\nu D/(l_\lambda l_\nu)\quad(\lambda\ne\nu),&&\\
 \operatorname{Var}_t(Y_\lambda)
   &=(1+t^2b_\lambda^2D_\lambda)/l_\lambda.&&
\end{align*}
All are polynomials. The off-diagonal denominators divide $D$;
the remaining two divisibilities are exactly
\eqref{eq:fibrepolecertificate}. Thus no pole is inferred to
cancel merely from the rational module's rank.

Finally integrate the zero-block markers conditionally. Its
mean $tC_0B_0(u)$ and covariance $C_0$ are polynomial in $t,u$.
Every original marker becomes a polynomial in $t,u,Y_\lambda$.
The polynomial Gaussian means and covariances above make its
period polynomial as well. This proves safety.

For the stated bound, if the zero block has nilpotency index
$r_0$, its mean has $t$ degree at most $r_0\le s$ and covariance
at most $r_0-1$. Since $\deg D+r_0\le s$, the joint means
above have degree at most $s+1$ and covariances at most $s$.
If the zero block is absent, the same bounds hold directly.
Conditional Wick expansion of a marker of degree $q$ has
coefficient degree at most $sq$ in $t$ and remaining variable
degree at most $bq$. Its Gaussian expectation therefore has
degree at most $sq+(s+1)bq$, proving
\eqref{eq:fibrecutoff}.
\end{proof}

For a general matrix $A$, the rationality conclusion and the
identities remain valid. At each nonzero eigenvalue, purity forces its
largest Jordan block to equal the full algebraic multiplicity,
because that is the maximum possible pole order of $K$.
Thus there can be only one Jordan block there. These identities alone
do not justify marker-pole cancellation for a defective nonzero block.
The proof above establishes that cancellation for semisimple nonzero
blocks and permits an unrestricted nilpotent zero block.
